\documentclass[11pt]{article}
\usepackage[a4paper,margin=24mm]{geometry}
\usepackage{amsmath,amssymb,amsthm,booktabs,array}
\usepackage[T1]{fontenc}
\usepackage{lmodern,microtype}
\usepackage[colorlinks=true,linkcolor=blue,citecolor=blue,urlcolor=blue]{hyperref}
\usepackage{fancyhdr}
\pagestyle{fancy}\fancyhf{}\fancyhead[L]{Two-component mollification of $\xi'$}
\fancyhead[R]{Expert-audit draft}\fancyfoot[C]{\thepage}
\setlength{\headheight}{14pt}
\newtheorem{theorem}{Theorem}
\newtheorem{lemma}[theorem]{Lemma}
\newtheorem{proposition}[theorem]{Proposition}
\newcommand{\ii}{\mathrm i}
\newcommand{\dd}{\,\mathrm d}
\newcommand{\Rea}{\operatorname{Re}}
\newcommand{\Ima}{\operatorname{Im}}
\newcommand{\Qop}{q(-\partial_A)q(-\partial_B)}
\title{An explicit two-component mollifier for zeros of\texorpdfstring{ $\xi'$}{ xi-prime} on the critical line}
\author{}
\date{27 October 2026 \quad---\quad Draft for expert audit}
\hypersetup{pdftitle={An explicit two-component mollifier for zeros of xi-prime on the critical line},pdfauthor={},pdfsubject={Proof draft for expert audit; multiplicity-counted zeros}}
\begin{document}
\thispagestyle{plain}
\begin{center}
{\Large An explicit two-component mollifier for zeros of $\xi'$\\[3pt]
on the critical line}\par\vspace{8pt}
27 October 2026 \quad---\quad Draft for expert audit
\end{center}
\begin{abstract}
We give a proof draft for a lower bound of $88.45\%$ for the proportion of
zeros of the derivative of the Riemann xi-function on the critical line,
counted with multiplicity. The construction combines the argument-variation
method used by Wu with an unrestricted Conrey--Feng mollifier of fixed length
$T^{4/7-10^{-6}}$. Only the convolution $\mu*\Lambda*\Lambda$ is added to the
Conrey component. All polynomials are specified by exact finite decimals.
We give the specialized leading-pole functional and an interval certificate
for its value. The analytic transplantation is presented for specialist
verification; the finite certificate does not verify that transplantation.
No assertion about simple zeros or publication priority is made.
\end{abstract}

\section{Statement and explicit data}

Put
\[
 \xi(s)=H(s)\zeta(s),\qquad
 H(s)=\tfrac12s(s-1)\pi^{-s/2}\Gamma(s/2).
\]
Let $N(T)$ count the nontrivial zeros of $\zeta$ with $0<\Ima\rho\le T$,
with multiplicity, and let $N_{\xi',c}(T)$ count zeros of $\xi'$ in this
height range on $\Rea s=1/2$, also with multiplicity.
The standard count for $\xi'$ is $N(T)+O(\log T)$; thus the denominator
below also gives its asymptotic proportion of all zeros.

\begin{theorem}[Claim submitted for verification]\label{thm:main}
The explicit construction below gives
\[
 \liminf_{T\to\infty}\frac{N_{\xi',c}(T)}{N(T)}>0.8845.
\]
\end{theorem}

This updates the multiplicity-counted quantity in Wu's Section~3
\cite{Wu}, not a simple-zero proportion. In particular, it is a different
quantity from the $0.86864$ simple-and-on-line bound stated in
\cite[Remark~7.1]{AF}. We fix, without further optimization,
\begin{equation}\label{eq:constants}
 R=1.04440033,\qquad \theta=\frac47-10^{-6},\qquad L=\log T,
 \qquad a=\frac12-\frac RL,\quad y=T^\theta.
\end{equation}
The operator and two profiles are
\begin{align}
q(x)={}&1-2.27043254x+0.304358115x^2+0.16793368x^3\notag\\
 &\hspace{7mm}+4.34152016x^4-5.3105844x^5+1.7701948x^6,\label{eq:q}\\
P_1(x)={}&-0.30528514x+2.50038954x^2-1.12177152x^3-0.07333288x^4,
 \label{eq:p1}\\
P_2(x)={}&1.89779072x+0.56594804x^2+0.07023055x^3.\label{eq:p2}
\end{align}
Every displayed decimal in \eqref{eq:constants}--\eqref{eq:p2} is an exact
rational number. In particular $q(0)=P_1(1)=1$ and $P_1(0)=P_2(0)=0$.
With $x_n=\log(y/n)/\log y$, define
\begin{equation}\label{eq:psi}
 \Psi(s)=\sum_{n\le y}\frac{1}{n^{s+R/L}}
 \left\{\mu(n)P_1(x_n)
       +\frac{(\mu*\Lambda*\Lambda)(n)}{(\log y)^2}P_2(x_n)\right\}.
\end{equation}
Here $*$ denotes Dirichlet convolution and $\Lambda$ is the von Mangoldt
function. There is \emph{no} factor $\mu^2(n)$ in the second component and
no extra factorial. Non-squarefree terms are retained. Write
\begin{equation}\label{eq:U}
 U(s)=q(-L^{-1}\partial_s)\zeta(s).
\end{equation}
The operator acts only on $\zeta$, never on $\Psi$.

\section{The counting interface}

The following formulation separates the exact counting function from its
asymptotic differential operator. It is the argument principle and
Littlewood reduction of \cite[Section~3, (18)--(23)]{Wu}, with the right
boundary normalization checked for \eqref{eq:psi}.

\begin{lemma}\label{lem:count}
For the fixed data above, suppose
\begin{equation}\label{eq:moment}
 \frac1T\int_T^{2T}|U(a+\ii t)\Psi(a+\ii t)|^2\dd t=M+o(1),
\end{equation}
and the same mollified second moments for each $L^{-k}\zeta^{(k)}$,
$0\le k\le6$, are bounded. Then
\begin{equation}\label{eq:count}
 \liminf\frac{N_{\xi',c}(T)}{N(T)}\ge1-\frac{\log M}{R}.
\end{equation}
\end{lemma}

\begin{proof}
Set $r(u)=q(1/2-u)$. Exact polynomial arithmetic gives
\begin{equation}\label{eq:r}
 r(u)=0.09491351+0.997010185u+2.089866u^2
                -2.29671034u^4+1.7701948u^6.
\end{equation}
If $r(u)=\sum r_ju^j$, use the exact auxiliary function
\begin{equation}\label{eq:eta}
 \eta(s)=\sum_{j=0}^6r_jL^{-j}\xi^{(j)}(s),\qquad V(s)=\eta(s)/H(s).
\end{equation}
On the critical line, $\xi^{(2j)}$ is real and $\xi^{(2j+1)}$ is purely
imaginary. Thus
\[
 \Ima\eta(1/2+\ii t)=\frac{0.997010185}{L}\Ima\xi'(1/2+\ii t).
\]
The nonzero coefficient is essential. The phase-counting argument of Wu
therefore applies directly, without identifying his polynomial parameters
with an unverified exact auxiliary function. Argument variation is taken
upward on the line, passing boundary zeros on their east side. Each net
advance of $\pi$ detects a zero of $\xi'$; zeros common to $\eta$ and the
line are included with multiplicity.

On the fixed window $T\le t\le2T$, Stirling's formula gives, uniformly in a
fixed strip around $1/2\le\Rea s\le3$,
\[
 \frac{H^{(j)}(s)}{H(s)L^j}=2^{-j}+O(L^{-1}).
\]
Since $r(1/2+z)=q(-z)$, it follows that
\[
 V-U=\sum_{k=0}^6O(L^{-1})L^{-k}\zeta^{(k)},\qquad
 \frac1T\int_T^{2T}|(V-U)\Psi|^2\dd t=O(L^{-2}).
\]
The coefficients in the first error may depend on $t$ but are uniformly
bounded as stated. Consequently the moment of $V\Psi$ is $M+o(1)$.

At $\Rea s=3$, absolute convergence and the profile endpoints show
\[
 \Psi_1(s)=\zeta(s)^{-1}+O(L^{-1}),\qquad \Psi_2(s)=O(L^{-2}),
 \qquad V(s)=\zeta(s)+O(L^{-1}).
\]
For the first estimate, bound the profile error by $O(\log n/\log y)$
and also include the factor $n^{-R/L}$ and the truncation tail. For the
second, use $\sum_n|\mu*\Lambda*\Lambda|(n)n^{-3}<\infty$.
Hence $V\Psi=1+O(L^{-1})$ on the right boundary. The logarithmic integral
there is $O(T/L)$; it is not necessary that $\Psi\to1$.

Let $n_V^+$ count zeros of $V$ in the right rectangle
$1/2\le\Rea s\le3$, $T\le\Ima s\le2T$, using the boundary convention
of Wu. Its argument principle and
$\Delta\arg H=\pi(N(2T)-N(T))+O(\log T)$ give
\begin{equation}\label{eq:phase}
 N_{\xi',c}(2T)-N_{\xi',c}(T)
 \ge N(2T)-N(T)-2n_V^++O(\log T).
\end{equation}
An equivalent safe convention overcounts boundary zeros on the right.
Heights hitting zeros are treated by limiting admissible heights.
Polynomial growth of the fixed-degree functions, together with their
nonzero right-boundary normalization, gives the horizontal argument
integrals $O(\log T)$ by the usual Jensen argument
\cite[Section~2, (6)--(8)]{Wu}. There are no poles of $H$ or $\zeta$ in
these high positive-height rectangles.

Every zero counted by $n_V^+$ is a zero of $V\Psi$ and is at horizontal
distance at least $R/L$ from $\Rea s=a$. Littlewood's formula then gives
\[
 2\pi\frac RL n_V^+
 \le\int_T^{2T}\log|V(a+\ii t)\Psi(a+\ii t)|\dd t
       +O(T/L)+O(\log T).
\]
Apply Jensen's inequality here, \emph{before} replacing $V$ by $U$.
No comparison of their negative logarithmic tails is used. Inserting the
second-moment estimate and \eqref{eq:phase}, and using
$N(2T)-N(T)\sim TL/(2\pi)$, yields the local version of
\eqref{eq:count}. Add finitely many dyadic windows $[T/2^J,T]$, with their
own fixed $L$ and mollifier. First let $T\to\infty$, then $J\to\infty$;
the omitted count is a fraction $2^{-J}+o(1)$. No estimate uniform in $J$
is required.
\end{proof}

\section{Shifted moment and the specialized kernel}

We use the third coefficient case of Pratt--Robles--Zaharescu--Zeindler
\cite[Theorem~4.1, p.~18]{PRZZ}. For fixed smooth height weight $\Phi$,
small complex shifts $\alpha,\beta=O(L^{-1})$, and the two coefficient
sequences obtained by expanding \eqref{eq:psi}, the theorem gives
\begin{align}\label{eq:PRZZ}
 &\int\zeta(\tfrac12+\alpha+\ii t)\zeta(\tfrac12+\beta-\ii t)
           \Psi(a+\ii t)\overline{\Psi(a+\ii t)}\Phi(t/T)\dd t\notag\\
 &\quad=\sum_{d,e\le y}\frac{a_da_e}{[d,e]}
        \frac{(d,e)^{\alpha+\beta}}{d^\alpha e^\beta}
 \int\!\left\{\zeta(1+\alpha+\beta)
 +\zeta(1-\alpha-\beta)
 \left(\frac{2\pi de}{t(d,e)^2}\right)^{\alpha+\beta}\right\}
 \Phi(t/T)\dd t+O(E),
\end{align}
where $\Psi(a+\ii t)=\sum a_nn^{-1/2-\ii t}$ and
$E\ll T^\varepsilon(y^{7/4}+y^{7/8}T^{1/2})$.
Here $(d,e)$ and $[d,e]$ denote gcd and lcm. The polynomial-profile
extension and finite linear combinations are included in the proof on
p.~18. Expanding the two components therefore includes both cross terms,
not just their diagonal moments. For our $\theta$, the two error exponents
before $\varepsilon$ are $1-1.75\cdot10^{-6}$ and
$1-0.875\cdot10^{-6}$. A sufficiently small fixed $\varepsilon$ absorbs
all fixed logarithmic factors and leaves $o(T)$.

The formula is uniform on fixed compact sets of the normalized complex
shifts $A=\alpha L$, $B=\beta L$. The apparent poles are combined before
differentiating. Cauchy's formula allows the fixed operator $\Qop$ to be
applied to the main term; the normalized derivatives cancel all powers
of $L$. This proves the bounded derivative moments in Lemma~\ref{lem:count}
as well. Nonnegative moments pass from smooth weights to the sharp window
by upper and lower approximations, taking the height limit before the
width of the endpoint smoothing tends to zero. We allow a slightly
enlarged fixed height interval for the upper approximation; a finite
partition and comparable rescaling reduce it to the published support
convention. The fixed positive length margin absorbs these changes.

Here is a finite integral expression for the limiting main term. It
specifies precisely the expression to be checked against the arithmetic
asymptotic. Label the Conrey component by $0$ and the Feng component by
$2$, and put $F_0=P_1$, $F_2=P_2$. For Mellin variables $s,t$, set
$X=A+s$, $Y=B+t$, $h=s+t$. The four leading kernels are
\begin{align}
 K_{00}&=XY/h,\notag\\
 K_{20}&=XY/h^3-2Y/h^2+Y/(Xh),\qquad
 K_{02}=XY/h^3-2X/h^2+X/(Yh),\label{eq:kernels}\\
 K_{22}&=7XY/h^5-6(X+Y)/h^4
       +(X/Y+Y/X+8)/h^3\notag\\
 &\hspace{13mm}-2(X^{-1}+Y^{-1})/h^2+(XYh)^{-1}.\notag
\end{align}
In particular the inverse powers of $X,Y$ must not be discarded.

For $0\le w\le\theta$, define
\begin{align}
 D_1(A,w;F)&=F'(w/\theta)/\theta+AF(w/\theta),\notag\\
 D_0(A,w;F)&=F(w/\theta),\label{eq:D}\\
 D_{-1}(A,w;F)&=\int_0^we^{-Az}F((w-z)/\theta)\dd z.\notag
\end{align}
Extend linearly the following rule on every monomial in
\eqref{eq:kernels}:
\begin{equation}\label{eq:H}
 cX^eY^fh^{-n}\ \longmapsto\
 \frac{c}{\theta^{\ell+m}}
 \int_0^\theta\frac{v^{n-1}}{(n-1)!}
 D_e(A,\theta-v;F_\ell)D_f(B,\theta-v;F_m)\dd v.
\end{equation}
Its value on $K_{\ell m}$ is denoted $H_{\ell m}(A,B)$. Only
$e,f\in\{-1,0,1\}$ occur. Define
\begin{align}\label{eq:I}
 I(A,B)=\sum_{\ell,m\in\{0,2\}}
 \frac{H_{\ell m}(A,B)-e^{-A-B}H_{\ell m}(-B,-A)}{A+B},\notag\\
 M=\left.\Qop I(A,B)\right|_{A=B=-R}.
\end{align}
The ordered cross terms are both present. The dual term swaps the two
shifts; $q$ differentiates the entire quotient in \eqref{eq:I}.

\begin{proposition}\label{prop:mainterm}
For \eqref{eq:psi}, the constant in \eqref{eq:moment} is \eqref{eq:I}.
\end{proposition}

\begin{proof}[Derivation from the published shifted moment]
Insert the coefficient expansion in \eqref{eq:PRZZ} and use Mellin inversion
for each monomial of $F_\ell$. The leading zeta ratio has prefactor
$XY/h$. Its logarithmic derivative for one left auxiliary variable is
$h^{-1}-X^{-1}$, for one right variable it is $h^{-1}-Y^{-1}$, and for a
left--right pair it is $h^{-2}$. The distinct auxiliary variables represent
the copies of $\Lambda$. All other connected derivatives for these
multiaffine coefficients vanish. Matching $k$ left variables to right
variables has multiplicity $\binom{\ell}{k}(m)_k$. Thus the leading factor is
\[
 \frac{XY}{h}\sum_{k=0}^{\min(\ell,m)}
 \binom{\ell}{k}(m)_k h^{-2k}
 (h^{-1}-X^{-1})^{\ell-k}(h^{-1}-Y^{-1})^{m-k}.
\]
For $\ell,m\in\{0,2\}$ this expands to \eqref{eq:kernels}.
This is the leading-pole extraction in \cite[Section~6, (6.3), and
pp.~39--40]{PRZZ}. The arithmetic Euler factor and its derivatives are
treated by that paper's contour argument: after the logarithmic
normalization, the omitted terms contribute $O(T/L)$, not a conjectural
ratios error. See also the procedure on p.~45. Replacing those factors by
one without this error estimate would not suffice.

Use
$h^{-n}=\int_0^\infty e^{-hv}v^{n-1}\dd v/(n-1)!$ in the Mellin inversions.
The profiles have support $w\ge0$ and vanish at zero. Multiplication by
$A+s$, by $1$, and by $(A+s)^{-1}$ therefore gives respectively the three
operators \eqref{eq:D}, with no endpoint delta term. Their common support
restricts $v$ to $[0,\theta]$. The factors $(\log y)^{-\ell-m}$ give
$\theta^{-\ell-m}$ after rescaling, proving \eqref{eq:H}.
The second term of \eqref{eq:PRZZ} gives the dual shift
$(-B,-A)$ and the factor $e^{-A-B}$, with the minus sign from the pole of
$\zeta(1-\alpha-\beta)$. This gives \eqref{eq:I}; the contour estimates and
shift uniformity above justify applying $\Qop$.
\end{proof}

As a normalization check, deleting the Feng component reduces \eqref{eq:I} to
\[
 1+\theta^{-1}\int_0^1\!\int_0^1 e^{2Rv}
 [P_1'(x)q(v)+\theta P_1(x)(Rq(v)+q'(v))]^2\dd x\dd v.
\]
This check is not a substitute for the arithmetic identification in
Proposition~\ref{prop:mainterm}.

\section{Finite certificate and completion}

\begin{lemma}\label{lem:finite}
For the rational data \eqref{eq:constants}--\eqref{eq:p2}, the functional
\eqref{eq:I} satisfies the outward enclosure
\begin{equation}\label{eq:interval}
 1.12818163807203565629105
 <M<1.12818163807203565629124.
\end{equation}
Moreover,
\begin{equation}\label{eq:threshold}
 1.12820541062467513417945
 <e^{R(1-0.8845)}<1.12820541062467513417947.
\end{equation}
\end{lemma}

\begin{proof}[Finite verification and its error bound]
The attached interval evaluator computes \eqref{eq:H} without quadrature.
It forms normalized Taylor jets in both shifts through order
six. In $D_{-1}$, expand $e^{-Az}$ through degree $N=24$. For
$C_F=\sum_d|F_d|$, the error in its $a$th normalized shift derivative at
$A=\pm R$ is bounded by
\begin{equation}\label{eq:tail}
 \frac{3C_F\theta^{1+a}}{a!(1+a)}
       \frac{(R\theta)^{N+1}}{(N+1)!},\qquad 0\le a\le6,
\end{equation}
because $R\theta<1$ and $e^{R\theta}<3$. The corresponding full-derivative
bound is the first factor in \eqref{eq:tail}. The finite polynomial
integrals are evaluated by
\[
 \int_0^w z^j(w-z)^d\dd z=\frac{j!d!}{(j+d+1)!}w^{j+d+1},\qquad
 \int_0^\theta\frac{v^{n-1}(\theta-v)^d}{(n-1)!}\dd v
   =\frac{d!}{(n+d)!}\theta^{n+d}.
\]
For a product of two truncated factors, combine the two errors with their
full-factor bounds and the product of errors; contract with absolute
values of the rational kernel and operator coefficients. The summed
analytic error is less than $8.799\cdot10^{-23}$.
All remaining operations, including the dual exponential $e^{2R}$,
use outward-rounded \texttt{mpmath.iv} intervals at 55 decimal digits.
They yield a narrower interval than \eqref{eq:interval}. An independent
40-digit, degree-20 setting also passes the threshold. The wrapper
\texttt{verify.py} checks the stated enclosures, exact polynomial
constraints, and the four expanded kernels. This is a numerical
certificate using a specified interval library, not a formal verification
of that library or of Proposition~\ref{prop:mainterm}.
\end{proof}

Combining Lemma~\ref{lem:count}, Proposition~\ref{prop:mainterm}, and
Lemma~\ref{lem:finite} proves the proposed Theorem~\ref{thm:main}.
The strict gap in \eqref{eq:interval}--\eqref{eq:threshold} is greater than
$2.37725\cdot10^{-5}$ in moment units. Thus the conclusion uses a fixed
rational witness, not convergence or global optimality of a numerical
optimizer. Direct high-precision evaluation gives
$1-\log M/R=0.884520175542196\ldots$, but the claimed bound is $0.8845$.

\section*{Verification priorities for the referee}

Review should prioritize the phase-counting convention in \eqref{eq:phase},
the shift-uniform applicability of \eqref{eq:PRZZ} to \eqref{eq:psi}, and
Proposition~\ref{prop:mainterm}. The code establishes only
Lemma~\ref{lem:finite}. The separate audit guide gives reproduction
commands, review questions and the literature-search scope.

\begingroup\small
\begin{thebibliography}{9}
\bibitem{Wu}
X. Wu, \emph{Distinct zeros of the Riemann zeta-function},
arXiv:1206.3737v2 (2012), especially Section~3, pp.~6--8.
\url{https://arxiv.org/pdf/1206.3737v2}.
\bibitem{PRZZ}
K. Pratt, N. Robles, A. Zaharescu and D. Zeindler,
\emph{More than five-twelfths of the zeros of $\zeta$ are on the critical line},
arXiv:1802.10521v3, Theorem~4.1 and Sections~6--7.
\url{https://arxiv.org/pdf/1802.10521v3}.
\bibitem{AF}
L. Alp\"oge and R. Furman, \emph{More than two thirds of the zeta zeros
are simple and on the critical line}, arXiv:2608.13637v2 (2026), Remark~7.1.
\url{https://arxiv.org/html/2608.13637v2#S7.SS3}.
\end{thebibliography}
\endgroup
\end{document}
